\begin{abstract}
Depth-pruned language models are commonly continued-pretrained to restore
aggregate language-model quality, but lower perplexity does not by itself
identify which behaviors have returned. We report an observational case study
of OLMo-2 prune--regrow recovery: retain a pretrained prefix, append two fresh
blocks, continue pretraining, and follow likelihood and capability over the
available checkpoints. The evaluation combines held-out likelihood, two MMLU
interfaces, three closed-book QA tasks, and a broader likelihood suite. In the
single 16-layer keep14 run, perplexity reaches 10.561 after 200k steps---still
$1.428\times$ the intact model's 7.398---while answer-letter MMLU is .319 versus
.605. An intact 32-layer control remains near the base through its available
25k checkpoint, arguing against short-horizon corpus shift but not providing a
200k counterfactual. Complete-option scoring raises keep14 to .383; however, a
fully random 16-layer operating point obtains nearly the same content score
while remaining at chance on answer letters, exposing a fluency floor. The gap
also appears on PopQA, TriviaQA, and NQ-open. A coupled construction comparison,
non-contiguous ShortGPT-16, reaches .474 MMLU at 200k; it does not isolate block
selection, inherited-block count, final-block retention, or fresh-tail effects.
The contribution is therefore narrow: an OLMo recovery-path case study and a
control bundle showing why likelihood, target capability, scoring interface,
exact construction, and recovery budget should be reported separately. It does
not establish seed-stable dynamics or a general law of post-pruning recovery.
\end{abstract}
